\chapter{Visualization and Analysis Tools}
\label{ch:visualization}

A plot is a measurement of selected model quantities. Its meaning depends on
coordinate frames, units, scaling and what has been omitted. Visualization can
expose an inconsistency or suggest a hypothesis; an attractive trajectory alone
does not validate dynamics, establish causality or reveal a golfer's control
strategy.

\section{Geometry Before Animation}

For a planar serial chain, let $q_1$ measure the first link from the positive
$x$ axis and let subsequent $q_i$ be relative joint angles. Then the absolute
orientation of link $i$ is $\alpha_i=\sum_{j=1}^{i}q_j$, and
\begin{equation}
  \bm{p}_i=\bm{p}_{i-1}
       +L_i(\cos\alpha_i,\sin\alpha_i)^T,\qquad \bm{p}_0=\bm{0}.
\end{equation}
A distal angle drawn from a world-horizontal reference is $\alpha_i$, not
$q_i$. This distinction also changes angular velocities used in energy.

\begin{figure}[ht]
  \centering
  \begin{tikzpicture}[scale=2]
    \coordinate (O) at (0,0);
    \coordinate (A) at ({1.5*cos(30)},{1.5*sin(30)});
    \coordinate (B) at ({1.5*cos(30)+cos(10)},{1.5*sin(30)+sin(10)});
    \draw[gray,dashed] (O) -- (0.8,0);
    \draw[very thick,blue] (O) -- (A) node[midway,above left] {$L_1$};
    \draw[very thick,red] (A) -- (B) node[midway,above] {$L_2$};
    \draw[gray,dashed] (A) -- ++(30:0.8);
    \draw[->] (0.6,0) arc[start angle=0,end angle=30,radius=0.6];
    \node at (0.8,0.17) {$q_1$};
    \draw[->] ($(A)+(30:0.6)$) arc[start angle=30,end angle=10,radius=0.6];
    \node at ($(A)+(20:0.95)$) {$q_2$};
    \fill (O) circle (0.025); \fill (A) circle (0.025);
    \fill (B) circle (0.025);
  \end{tikzpicture}
  \caption{Relative Joint Angles: $q_1=30^\circ$, $q_2=-20^\circ$,
  and Distal Absolute Orientation $\alpha_2=10^\circ$}
  \label{fig:stick-figure}
\end{figure}

\begin{lstlisting}[language=Python, caption={Planar Joint Positions With Relative Angles}]
import numpy as np

def compute_joint_positions_2d(
    q: np.ndarray, lengths: np.ndarray,
) -> np.ndarray:
    """Return (n+1,2) positions, including a fixed origin."""
    angles = np.asarray(q, dtype=float)
    size = np.asarray(lengths, dtype=float)
    if angles.ndim != 1 or angles.size == 0:
        raise ValueError("q must be a nonempty vector")
    if size.shape != angles.shape or np.any(size <= 0):
        raise ValueError("lengths must match q and be positive")
    if not np.isfinite(angles).all() or not np.isfinite(size).all():
        raise ValueError("angles and lengths must be finite")
    absolute = np.cumsum(angles)
    links = size[:, None] * np.column_stack(
        (np.cos(absolute), np.sin(absolute))
    )
    return np.vstack((np.zeros(2), np.cumsum(links, axis=0)))
\end{lstlisting}

This computes positions, not an animation and not a three-dimensional
reconstruction. An animation additionally needs a time base, interpolation,
fixed axes and a renderer. A 3D chain requires the actual joint axes and
rigid transforms; changing viewing angle cannot recover missing depth from a
planar model. Plot COM locations and contact points separately from joints.
Preserve lengths and compare analytic or finite-difference endpoint velocities
with the declared Jacobian before interpreting motion.

\section{A Phase Projection Is Not the Full State}

The following helper preserves the caller's angle convention and returns
copies. It does not unwrap data. For a coupled multi-joint system, the pair
$(q_i,\dot q_i)$ is generally a projection: two visits to the same plotted
point can have different accelerations because other coordinates, actuator
states or inputs differ. Crossing curves need not violate uniqueness of the
full dynamics. Wrapping a revolute angle can introduce display jumps;
unwrapping can hide the physical periodicity. State the choice and use time
or event markers when comparing swings.

\begin{lstlisting}[language=Python, caption={An Explicit Joint-State Projection}]
def phase_portrait(
    q_history: np.ndarray, qd_history: np.ndarray,
    joint_idx: int = 0, label: str = "Joint 0",
) -> dict:
    """Project finite matching (N,n) histories without wrapping."""
    q = np.asarray(q_history, dtype=float)
    v = np.asarray(qd_history, dtype=float)
    if q.ndim != 2 or min(q.shape) < 1 or v.shape != q.shape:
        raise ValueError("histories must have matching nonempty (N,n) shape")
    if not np.isfinite(q).all() or not np.isfinite(v).all():
        raise ValueError("histories must be finite")
    if (isinstance(joint_idx, bool)
            or not isinstance(joint_idx, (int, np.integer))
            or not 0 <= joint_idx < q.shape[1]):
        raise ValueError("joint index is outside the history")
    return {
        "x": q[:, joint_idx].copy(),
        "y": v[:, joint_idx].copy(),
        "label": label,
    }
\end{lstlisting}

\section{A Velocity Ellipsoid Needs a Metric}

For a fixed configuration, $\bm{v}_{\mathrm{task}}=J\dot{\bm{q}}$ is
an instantaneous kinematic map. Choose positive rate scales $s_i$ and
assume $\|\bm{z}\|_2\leq1$ with $\dot{\bm{q}}=D\bm{z}$,
$D=\operatorname{diag}(s_i)$. The attainable set under this assumption is
\begin{equation}
  \mathcal{E}=\{JD\bm{z}:\|\bm{z}\|_2\leq1\}.
\end{equation}
If $A=JD=U\Sigma V^T$, its principal directions are columns of $U$
and its semiaxis lengths are the singular values of $A$, or the square roots
of eigenvalues of $AA^T$. The matrix $AA^T$ describes the ellipsoid's shape;
it is not itself the set \cite{ModernRoboticsManipulabilityReview}.

\begin{lstlisting}[language=Python, caption={The Image of a Scaled Joint-Rate Ball}]
def velocity_ellipsoid(
    jacobian: np.ndarray, rate_scales: np.ndarray,
) -> tuple:
    """Return full task axes, radii (including zeros), and rank."""
    matrix = np.asarray(jacobian, dtype=float)
    scales = np.asarray(rate_scales, dtype=float)
    if matrix.ndim != 2 or min(matrix.shape) < 1:
        raise ValueError("Jacobian must be a nonempty matrix")
    if scales.shape != (matrix.shape[1],) or np.any(scales <= 0):
        raise ValueError("one positive scale is required per joint")
    if not np.isfinite(matrix).all() or not np.isfinite(scales).all():
        raise ValueError("Jacobian and scales must be finite")
    mapped = matrix * scales
    if not np.isfinite(mapped).all():
        raise ValueError("scaled Jacobian overflow")
    axes, singular, _ = np.linalg.svd(mapped, full_matrices=True)
    radii = np.zeros(matrix.shape[0])
    radii[:singular.size] = singular
    return axes, radii, int(np.linalg.matrix_rank(mapped))
\end{lstlisting}

At a singular configuration the image collapses into a lower-dimensional set;
do not invert $AA^T$ blindly. Repeated singular values make individual axis
directions nonunique. The returned rank uses NumPy's numerical tolerance,
so report scales and tolerance when interpreting a nearly singular result.
Independent bounds $|\dot q_i|\leq s_i$ form a box, whose linear image is
generally a zonotope, not this Euclidean-ball ellipsoid. The ball is contained
in that box and encodes an additional combined-rate budget.

Translational and angular task velocities have different units. Analyze them
separately or define a task metric, such as an explicitly justified length
scale, before comparing their magnitudes. A large kinematic semiaxis means
large instantaneous task velocity under the chosen rate budget; it does not
mean low torque, low energy, dynamic controllability or physiological ease.
For golf, clubhead translation, face orientation and strike direction can
compete for the same joint-rate resources.

\section{Energy Accounting and Transfer}

For rigid segments, include both COM translation and rotation:
\begin{equation}
  T=\sum_i\left[
    \frac12m_i\|\bm{v}_{C_i}\|^2+
    \frac12\bm{\omega}_i^T I_{C_i}^{\,W}\bm{\omega}_i
  \right]=\frac12\bm{v}^TM(\bm{q})\bm{v},
  \qquad U_g=\sum_i m_i g h_{C_i}.
\end{equation}
The inertia and angular velocity in each term use the same frame; the
superscript $W$ denotes world coordinates. For planar relative joints,
$\omega_i=\sum_{j\leq i}\dot q_j$, not $\dot q_i$ alone.
COM velocities also depend on upstream motion. Independent terms
$\tfrac12 I_i\dot q_i^2$ therefore miss translational energy and coupling
unless a specifically derived model justifies that decomposition.

With ideal fixed supports and conservative gravity, add any conservative
elastic potential to $U$. The work balance becomes
\begin{equation}
  \frac{d}{dt}(T+U)
    =\bm{v}^T\bm{\tau}_{\mathrm{act}}
       +P_{\mathrm{ext}}-P_{\mathrm{diss}},
\end{equation}
where actuator torque is the generalized force conjugate to $\bm{v}$,
external power includes moving supports and applied wrenches, and
$P_{\mathrm{diss}}\geq0$ represents the declared dissipative elements.
Ideal constraints do no net power only under the corresponding stationary,
workless constraint assumptions. Impacts require an energy jump balance.
Zero muscle torque alone does not imply conservation: damping, moving
contacts or other actuators may remain. Conservative passive elasticity can
exchange energy with $T+U_g$ while conserving the augmented $T+U$.

An internal joint torque pair has net power
$\bm{\tau}\cdot(\bm{\omega}_{\mathrm{distal}}
-\bm{\omega}_{\mathrm{proximal}})$ under the stated sign convention.
Intersegmental force and moment power depends on the point and frame used;
include force work at a moving joint when drawing segment energy flows.
An increase in a distal segment's energy alone does not prove a proximal
muscle supplied it. Check total work balance first, then define the segment
boundaries and all crossing powers. Chapter~\ref{ch:golf-capstone} applies
the canonical model to signed acceleration contributions and prescribed-motion
torque requirements; those quantities answer different questions from energy
transfer or neural coordination.

\section*{Exercises}
\begin{enumerate}
  \item Evaluate a two-link chain at $(q_1,q_2)=(0,\pi/2)$ and
  $(0,\pi)$. Compare a finite-difference endpoint Jacobian with its analytic
  derivative and identify a singular pose.
  \item Change one joint-rate scale and plot the resulting ellipsoid. Compare
  it with the mapped corners of the independent joint-rate box.
  \item For two uniform rods, calculate COM translational and rotational
  energies from geometry and compare their sum with
  $\tfrac12\bm{v}^TM\bm{v}$. Explain why using relative joint rates as
  absolute segment rates fails.
\end{enumerate}
